% BEGIN PROFESSOR FEEDBACK PF09: Method - Data acquisition
\section{Data acquisition}
\label{sec:data_acquisition}

Studying coordinated malicious behaviour on-chain needs two things at once:
transaction data, and reliable labels that say which addresses are malicious.
This combination is scarce. Transactions are public, but the intent behind an
address can only be established off-chain, through abuse reports, exchange
records, or investigations, and such labels are rare and unevenly spread across
the network \citep{shi2025realcats}. To work around this, the reported
evaluation uses a main Ethereum dataset and the category-labelled Real-CATS
dataset. Separate synthetic templates check the correctness of subgraph
counting. The pipeline also supports BLTE as an optional secondary
check, but no BLTE result is used to answer the research questions in this
thesis. The origin and prior use of the two published datasets, BLTE and
Real-CATS, were already covered in Section~\ref{sec:relwork_datasets}. Here each
real-world dataset is described by its source, its contents, and its
suitability for this thesis. Validation templates and the distinction between
predictive tests and network null models are described separately below.

\subsection{Real-world datasets}

\paragraph{Hugging Face (HF) Ethereum Fraud Dataset.}
\label{sec:dataset_hf}
\textit{Source.} The dataset is published on the Hugging Face Hub
\citep{fesevu2025ethfraud}. This thesis uses its full-history release, which
pairs a complete transaction table with a set of labelled addresses, rather than
one of the smaller monthly extracts.

\textit{Contents.} The dataset has two parts. The first is a transaction table
of 380{,}566{,}661 transfers recorded between August 2015 and August 2025. Each
transfer carries:
\begin{itemize}
  \item the sending and receiving address, as hashed identifiers,
  \item the timestamp of the transfer,
  \item the transferred value,
  \item the transaction hash and block number,
  \item a flag marking contract-creation transactions, and
  \item the transaction fee.
\end{itemize}
The second part is a label table of 114{,}739 addresses, of which 60{,}086 are
benign and 54{,}653 fraudulent. Each labelled address also states whether it is
a contract or an ordinary account, and carries a short text description of the
reported activity.

\textit{Suitability.} This is the main dataset of the thesis, and it carries the
evaluation for RQ1 and RQ2. It is the only dataset used here that
combines full timestamped transaction histories with a large and roughly
balanced set of malicious and benign labels. The timestamps matter because
temporal motifs depend on the exact timing and order of transfers, and the size
of the label set is large enough to draw more than one cohort of addresses while
still keeping an untouched test set aside.

\paragraph{BLTE.}
\label{sec:dataset_blte}
\textit{Source.} BLTE was introduced by \citet{alemari2021labeled} and is built
on the Ethereum Classic chain. Its origin and the results reported on it by
earlier work were described in Section~\ref{sec:relwork_datasets}. Here it is
used as obtained from the public release.

\textit{Contents.} The dataset contains 71{,}250 transactions from October 2017
to December 2019, labelled one transaction at a time as normal or abnormal.
There are 57{,}000 normal and 14{,}250 abnormal transactions, so about one
transaction in five is abnormal. Each record provides:
\begin{itemize}
  \item the transaction hash and nonce,
  \item the sending and receiving address,
  \item the transferred value, the gas, and the gas price,
  \item the block timestamp, block number, and block hash, and
  \item four labelling fields that mark whether the sender or the receiver is a
        known scam, and of which category.
\end{itemize}
A set of pre-computed per-address features is distributed alongside the raw
transactions, but this thesis does not use them, for the reason given below.

\textit{Suitability.} BLTE can be used as an optional secondary check of whether
the main finding also appears in another Ethereum-based dataset. It is less
directly comparable with the main dataset because its labels apply to
transactions rather than addresses and because only 169 of its 73{,}034 active
addresses are marked as scams. A BLTE run is therefore kept separate from the
main evaluation and uses class weighting to account for this imbalance. The
pre-computed features are left aside so that the same behavioural features can
be computed as for the other datasets. No completed BLTE result is included in
the Results chapter, and BLTE is not needed to answer any of the four research
questions.

\paragraph{Real-CATS.}
\label{sec:dataset_realcats}
\textit{Source.} Real-CATS was introduced by \citet{shi2025realcats}. Its design
and the sources of its labels were described in
Section~\ref{sec:relwork_datasets}. This thesis uses its Ethereum portion
together with the detailed transaction records released alongside it.

\textit{Contents.} The Ethereum portion covers 12{,}572 criminal and 16{,}019
benign addresses. For each address the dataset provides its transaction history,
split into five separate layers: ordinary transfers, internal transfers, and
ERC-20, ERC-721, and ERC-1155 token transfers. What sets Real-CATS apart from
the other datasets is that its criminal addresses carry a specific category,
rather than a single malicious flag. The categories are:
\begin{itemize}
  \item Phishing Scam and Impersonation Scam,
  \item Contract Exploit Scam and Hack Scam,
  \item Blackmail Scam and Fake Returns Scam,
  \item Rug Pull Scam, NFT Airdrop Scam, and Pigbutchering Scam.
\end{itemize}

\textit{Suitability.} Real-CATS is the dataset used for RQ3 and RQ4. Its detailed
categories are what make both questions possible, because they allow the motif
fingerprints of different reported malicious behaviours to be compared directly,
rather than only a single malicious-versus-benign split. The labels do not state
whether an address belongs to a coordinated campaign, so the experiment tests
category-level fingerprint differences rather than independently verified
coordination. The labels are also
independent of the main dataset, since they come from public abuse reports
instead of the process behind the HuggingFace release. The categories are very
unevenly filled, with phishing and impersonation together accounting for most of
the criminal addresses, so RQ3 and RQ4 are limited to the categories that have
enough addresses to support a fair comparison and treat the smaller ones as
descriptive only. The catch-all category \emph{Other} is excluded as well, even
though it is large enough to qualify, because it collects unrelated reports and
therefore does not describe one kind of behaviour.

\begin{table}[ht]
\centering
\caption{The real-world datasets and their roles. The malicious share is
computed over addresses. Synthetic validation templates are described
separately in Section~\ref{sec:dataset_synthetic}.}
\label{tab:dataset_comparison}
\small
\setlength{\tabcolsep}{4pt}
\begin{tabular}{@{}p{3.1cm} p{3.1cm} p{3.0cm} p{3.2cm}@{}}
\toprule
\textbf{Property} & \textbf{HuggingFace} & \textbf{BLTE} & \textbf{Real-CATS} \\
\midrule
Chain              & Ethereum      & Ethereum Classic & Ethereum \\
Transactions       & 380.6 million & 71{,}250         & per-address histories \\
Label level        & address       & transaction      & address \\
Address coverage   & 114{,}739 labelled & 73{,}034 active & 28{,}591 labelled \\
Malicious share    & 47.6\,\%      & 0.23\,\%         & 44.0\,\% \\
Date range         & 2015--2025    & 2017--2019       & per release \\
Category analysis  & no            & no               & 9 retained categories \\
Role               & RQ1, RQ2      & optional check   & RQ3, RQ4 and binary check \\
\bottomrule
\end{tabular}
\end{table}

\subsection{Validation templates and null models}
\label{sec:validation_and_null_models}

\paragraph{Synthetic validation templates.}
\label{sec:dataset_synthetic}
The validation data is generated in code. Each template contains exactly three
directed transfers on two or three nodes, with known directions and timestamps
and no additional background edges. The templates systematically cover the
patterns accepted by the counter. They do not simulate complete fraud campaigns
or provide malicious-address labels. A three-node structure may resemble one
part of a plausible scam, such as collecting and forwarding funds, but it
remains a candidate motif before any statistical comparison.

Raphtory returns an ordered vector of 40 local motif-role counts for each
address. A position means one component of that vector, not one of 40 global
patterns. The templates check all 24 star-centre roles, eight two-node roles,
and eight triangle roles. The difference from the 36 global patterns of
\citet{paranjape2017motifs} is explained in
Section~\ref{sec:motifcountingandretainedroles}. The templates also check
transfers exactly at a window boundary and transfers with equal timestamps.
Their purpose is to confirm that subgraph counting recognises a known two- or
three-node structure and assigns the expected local role counts. They are
never used to measure detection performance or compared with real-world
datasets as examples of fraud.

\paragraph{Network null models.}
A validation template has a known answer for the counter. A network null model
serves a different purpose: it generates random reference networks against which
the observed frequency of a candidate pattern can be tested
\citep{milo2002network}. For example, shuffling timestamps while retaining the
observed directed transfers would test their temporal arrangement conditional
on that fixed structure. The choice of preserved properties determines the
claim that such a comparison can support.

The completed experiments do not include a random-network null model, so no
individual candidate is claimed to be statistically overrepresented. The
permutation and bootstrap procedures in
Section~\ref{sec:evaluation_protocol} instead assess predictive differences and
category separability. The alternative orderings of equal-timestamp transfers
are sensitivity checks, not random-network references.
% END PROFESSOR FEEDBACK PF09

